\documentclass[11pt,reqno]{amsart}

\usepackage[a4paper,margin=1in]{geometry}
\usepackage{amsmath,amsthm}
\usepackage{mathtools}
\usepackage{latexsym}
\usepackage{amsfonts}
\usepackage{amssymb}
\usepackage[dvipsnames]{xcolor}
\usepackage{graphicx}
\usepackage{float}
\usepackage{graphbox}
\usepackage{textcomp}
\usepackage[colorlinks=true,linkcolor=blue,urlcolor=blue,citecolor=red]{hyperref}
\usepackage{array}
\usepackage{comment}
\usepackage{appendix}
\usepackage{enumitem}

\newtheorem{theorem}{Theorem}[section]
\newtheorem{proposition}[theorem]{Proposition}
\newtheorem{lemma}[theorem]{Lemma}
\newtheorem{corollary}[theorem]{Corollary}
\newtheorem{definition}[theorem]{Definition}
\newtheorem{question}[theorem]{Question}
\newtheorem{remark}[theorem]{Remark}
\newtheorem{example}[theorem]{Example}

\newcommand{\norm}[1]{\left\lVert#1\right\rVert}
\newcommand{\abs}[1]{\left\lvert#1\right\rvert}
\def\R{\mathbb{R}}
\def\C{\mathbb{C}}
\def\N{\mathbb{N}}
\DeclareMathOperator{\Tr}{Tr}
\DeclareMathOperator{\E}{\mathbb{E}}
\DeclareMathOperator{\Prob}{P}
\DeclareMathOperator{\diag}{diag}

% Draft provenance colors.
% Uncolored text is pre-existing draft material.
% `revderived`: assistant organization/paraphrase of pre-existing draft text.
% `revverified`: prose translated from verified proof-route evidence.
% `revnewmath`: genuinely new mathematical exposition added in this pass.
\definecolor{RevisionDerived}{RGB}{105,105,105}
\definecolor{RevisionVerified}{RGB}{0,105,120}
\definecolor{RevisionNewMath}{RGB}{0,45,150}
\newenvironment{revderived}{\begingroup\color{RevisionDerived}}{\endgroup}
\newenvironment{revverified}{\begingroup\color{RevisionVerified}}{\endgroup}
\newenvironment{revnewmath}{\begingroup\color{RevisionNewMath}}{\endgroup}
% Backwards-compatible aliases used by the included appendix body.
\newenvironment{revlatest}{\begin{revnewmath}}{\end{revnewmath}}
\newenvironment{revrecent}{\begin{revverified}}{\end{revverified}}
\newenvironment{revearlier}{\begin{revderived}}{\end{revderived}}

\title[Sharp upper tails for PPT moments]{Sharp upper tails for partial-transpose moments on the Hilbert--Schmidt sphere}
\author{}
\date{}

\begin{document}
\begin{revverified}
\begin{abstract}
We isolate the upper-bound half of a sharp large-deviation principle for the
\(k\)-th partial-transpose moment of a random induced state in the balanced
regime \(s(d)\sim \lambda d^2\).  The upper-tail speed is
\(d^{2+2/k}\), and the rate predicted by the one-column spike mechanism is
\(\lambda\varepsilon^{1/k}\).  The proof developed here explains why no
larger upper-tail probability can occur: after removing a background typical
set, the bad event is forced outside a spherical neighbourhood whose
isoperimetric cost is exactly
\(\exp\{-\lambda\varepsilon^{1/k}d^{2+2/k}+o(d^{2+2/k})\}\).
The article is organized as a conditional proof route; the older
\(d^2\)-scale concentration appendix is retained as a supplier for the
typical-set estimate, not as the sharp theorem.
\end{abstract}
\end{revverified}

\maketitle

\begin{revnewmath}
\begin{remark}[Draft provenance colors]
The colors are provenance labels, not decoration.  Uncolored text is inherited
from the pre-existing draft.  Gray text is assistant organization or paraphrase
of that draft.  Teal text is translated from proof-route statements.  Dark blue text is new mathematical exposition added in
this pass.  The color environments are local and can be redefined to no-ops
for a submission version.
\end{remark}
\end{revnewmath}

\begin{revnewmath}
\section{Scope and relation with the PPT threshold}
\label{sec:scope-aubrun}

Aubrun's theorem on partial transposition of random induced states is a
quantitative threshold result.  In particular, the statement is not merely
that the balanced PPT threshold is the constant \(4\): the theorem also gives
high-probability estimates, with explicit inequalities controlling the two
sides of the threshold.  The present paper should therefore not be read as a
replacement proof of Aubrun's threshold theorem unless those probability
bounds are restated and matched in the same normalization.

The target here is narrower.  We study a fixed moment order \(k\ge3\), the
balanced environment \(s(d)\sim\lambda d^2\), and the logarithmic upper tail
of the partial-transpose moment at speed
\[
  v_{k,d}=d^{2+2/k}.
\]
The mechanism is a large-deviation mechanism for a moment functional: a spike
of spherical radius \(d^{-1+1/k}\) creates a moment excess of order one, and
the ambient Hilbert--Schmidt sphere charges such a displacement with cost
\(\exp\{-\lambda a\,d^{2+2/k}+o(d^{2+2/k})\}\).  This is compatible with the
PPT-threshold picture, but it is a different assertion from Aubrun's
quantified threshold bounds.
\end{revnewmath}

\begin{revverified}
\section{Sharp logarithmic upper tail}
\label{sec:sharp-log-upper-tail}

Fix an integer \(k\ge3\), a level \(\varepsilon>0\), and a balanced
environment sequence \(s(d)\) with \(s(d)/d^2\to\lambda>0\).  Let \(X_d\) be
uniform on the Hilbert--Schmidt unit sphere in \(M_{d^2,s(d)}(\C)\), and let
\[
  p_d^\uparrow(\varepsilon)
  =
  \Prob\left\{
    \left|
      m_{k,d}(X_d)-\E m_{k,d}(X_d)
    \right|\ge \varepsilon
  \right\}
\]
denote the absolute-deviation upper-tail probability for the normalized
background moment used in the formal Appendix B development.  The speed and
candidate rate are
\[
  v_{k,d}=d^{2+2/k},
  \qquad
  I_k(\varepsilon)=\lambda\varepsilon^{1/k}.
\]

\begin{theorem}[Formalized sharp upper endpoint]
\label{thm:sharp-upper-endpoint}
Assume the following inputs for the concrete spherical matrix model:
\begin{enumerate}[leftmargin=*]
\item the full spherical isoperimetric inequality in the ambient real
  dimension \(2d^2s(d)\);
\item the dimension identifications between the matrix model and that real
  sphere;
\item eventual positive mass of the one-sided favourable upper-tail event;
\item eventual half-mass of the background typical set, obtained from the
  moment and operator-norm bad-set estimates;
\item the deterministic mixed-word expansion bound in the sharp spherical
  neighbourhood, together with the corresponding scalar term limits.
\end{enumerate}
Then for every \(\eta>0\), eventually in \(d\),
\[
  \frac{\log p_d^\uparrow(\varepsilon)}{v_{k,d}}
  \le
  -I_k(\varepsilon)+\eta .
\]
\end{theorem}
\end{revverified}

\begin{revnewmath}
\begin{theorem}[Main upper large-deviation bound]
\label{thm:main-upper-large-deviation}
Let \(k\ge3\), \(\varepsilon>0\), and \(s(d)/d^2\to\lambda>0\).  Suppose that
the five input estimates listed in Theorem~\ref{thm:sharp-upper-endpoint}
hold for the balanced spherical model.  Then
\[
  \limsup_{d\to\infty}
  \frac{1}{d^{2+2/k}}
  \log
  \Prob\left\{
    \left|
      m_{k,d}(X_d)-\E m_{k,d}(X_d)
    \right|\ge \varepsilon
  \right\}
  \le
  -\lambda\varepsilon^{1/k}.
\]
\end{theorem}

\begin{proof}
Apply Theorem~\ref{thm:sharp-upper-endpoint} with an arbitrary
\(\eta>0\), take the \(\limsup\), and then let \(\eta\downarrow0\).  The
content is therefore the formalized endpoint itself: all dimensional
normalizations, local expansion estimates, and spherical-isoperimetric
constants have already been reduced to the concrete model.
\end{proof}

\begin{proof}[Proof route]
Choose a slack parameter and set the spherical radius so that its \(k\)-th
power is just below \(\varepsilon\).  On the background typical set, the local
expansion of the \(k\)-th moment separates into the pure radial spike term,
the background term, and mixed words.  The deterministic mixed-word estimate
and its scalar limits make the mixed contribution smaller than the available
slack, so every point in the upper-tail event must lie outside the sharp
spherical neighbourhood of the typical set.

The typical set has probability at least \(1/2\).  Spherical isoperimetry
therefore bounds the probability of the complement of its radius-\(r_d\)
neighbourhood by
\[
  \exp\left(-\frac{(2d^2s(d)-1)r_d^2}{2}\right).
\]
The radius is chosen at the spike scale, and the scalar identities reduce the
exponent to
\[
  -\lambda\,a + o(1)
\]
after division by \(d^{2+2/k}\), where \(a^k\) is the admissible moment
increase.  Letting the slack tend to zero gives \(a\to\varepsilon^{1/k}\),
which yields the claimed \(-\lambda\varepsilon^{1/k}\) upper rate.
\end{proof}

\begin{lemma}[Radius-to-rate computation]
\label{lem:radius-to-rate}
Let \(N_d=d^2\), let \(s(d)/N_d\to\lambda\), and let
\(v_{k,d}=d^{2+2/k}\).  For a fixed \(a\ge0\), set
\[
  r_d^2=\frac{a\,v_{k,d}}{N_d^2},
  \qquad
  n_d=2N_ds(d).
\]
Then
\[
  \frac{(n_d-1)r_d^2}{2v_{k,d}}\longrightarrow \lambda a .
\]
\end{lemma}

\begin{proof}
Since \(r_d^2=a v_{k,d}/N_d^2\),
\[
  \frac{(n_d-1)r_d^2}{2v_{k,d}}
  =
  a\left(\frac{s(d)}{N_d}-\frac{1}{2N_d^2}\right).
\]
The first term tends to \(\lambda\), and the second tends to \(0\).
\end{proof}

\begin{theorem}[Sharp exponent, conditional on the matching lower spike]
\label{thm:sharp-optimal-exponent}
Assume, in addition to the hypotheses of
Theorem~\ref{thm:sharp-upper-endpoint}, the one-column lower spike
construction: for every \(\eta>0\), eventually in \(d\),
\[
  -I_k(\varepsilon)-\eta
  \le
  \frac{\log p_d^\uparrow(\varepsilon)}{v_{k,d}} .
\]
Then
\[
  \lim_{d\to\infty}
  \frac{\log p_d^\uparrow(\varepsilon)}{d^{2+2/k}}
  =
  -\lambda\varepsilon^{1/k}.
\]
\end{theorem}

\begin{proof}
The lower spike gives
\(\liminf_d v_{k,d}^{-1}\log p_d^\uparrow(\varepsilon)
\ge -I_k(\varepsilon)\).  The sharp upper endpoint gives
\(\limsup_d v_{k,d}^{-1}\log p_d^\uparrow(\varepsilon)
\le -I_k(\varepsilon)\).  The two inequalities force the limit and identify
the optimal exponent.
\end{proof}

\begin{remark}[What is stronger than the concentration appendix]
The concentration theorem later in this file controls ordinary fluctuations
around the mean on an exponential scale supplied by the ambient sphere.  The
sharp endpoint above is a large-deviation statement at the smaller spike
speed \(d^{2+2/k}\), and it identifies the exact constant in the exponent.
Thus the concentration argument is only one supplier for the typical
background set; it is not the final optimality theorem.
\end{remark}
\end{revnewmath}

\begin{revverified}
\section{Proof bookkeeping}
\label{sec:proof-bookkeeping}

The proof separates the probabilistic argument into an upper inequality and a
matching one-column lower construction.  Both sides are normalized at the same
speed and rate:
\[
  v_d=d^{2+2/k},
  \qquad
  I(\varepsilon)=\lambda\varepsilon^{1/k}.
\]
The bookkeeping point is that the upper argument is an optimal-exponent
argument, not merely an ordinary concentration estimate: the concentration
appendix supplies typical background control, while the sharp speed comes from
the spike displacement and the cap cost on the Hilbert--Schmidt sphere.
\end{revverified}

\begin{revderived}
\section{Introduction}

Let \(X\) be uniformly distributed on the Hilbert--Schmidt unit sphere in
\(M_{d^2,s}(\C)\).  The matrix \(XX^*\) is the normalized induced random
state obtained from a Gaussian matrix after radial projection, and
\((XX^*)^\Gamma\) denotes its partial transpose with respect to the tensor
decomposition \(M_{d^2}(\C)=M_d(\C)\otimes M_d(\C)\).  The quantities studied
in this article are the moments
\[
  f_k(X)=\Tr\left(((XX^*)^\Gamma)^k\right),
  \qquad k\ge1.
\]

The balanced regime is \(s\asymp d^2\).  This is the natural threshold scale
for partial transposition of random induced states: the matrix dimension is
\(d^2\), the number of environmental columns is comparable to \(d^2\), and the
real dimension of the Hilbert--Schmidt sphere is
\[
  2d^2s-1\asymp d^4.
\]
This large ambient dimension is the source of the concentration estimate
below.

There is, however, a useful obstruction.  The map \(X\mapsto f_k(X)\) is not
treated as globally Lipschitz at the optimal scale.  A global Lipschitz bound
would be too crude, because it would have to control rare configurations with
large operator norm.  Instead, the proof constructs a high-probability set
\(\Omega\) on which the relevant operator norms are of order \(d^{-2}\).  On
this set the trace-power perturbation estimate gives the correct local
Lipschitz constant.  A localized form of Levy's lemma then converts the local
Lipschitz estimate and the smallness of \(\Omega^c\) into concentration on the
whole sphere.

\section{Model and main result}

The ambient space is the Hilbert--Schmidt unit sphere
\[
  S_{HS}(d^2,s)
  =
  \{X\in M_{d^2,s}(\C):\norm X_2=1\},
\]
with its uniform surface probability measure.  The observable is
\[
  f_k(X)=\Tr\left(((XX^*)^\Gamma)^k\right),
\]
where \(\Gamma\) denotes partial transposition on the second tensor factor of
\(M_d(\C)\otimes M_d(\C)\).

\begin{theorem}[Main concentration estimate]
\label{thm:main-concentration}
Fix \(k\ge1\) and \(\lambda>0\).  There are constants \(c,C>0\), depending
only on \(k\) and \(\lambda\), such that the following holds for all
sufficiently large \(d\).  If
\[
  \frac{\lambda}{2}d^2\le s\le 2\lambda d^2
\]
and \(X\) is uniform on \(S_{HS}(d^2,s)\), then for every \(\varepsilon>0\),
\[
  \Prob\left\{
    \left|
      f_k(X)-\E f_k(X)
    \right|
    \ge
    \varepsilon d^{-(2k-2)}
  \right\}
  \le
  C e^{-c d^2}
  +
  C\exp\left(-c\frac{d^4\varepsilon^2}{k^2}\right).
\]
\end{theorem}

The first term in the bound is the cost of discarding the bad set where the
operator-norm controls may fail.  The second term is the spherical
concentration term.  The exponent \(d^4\) is the real dimension of the sphere,
while the factor \(k^2\) comes from the local Lipschitz constant of the
degree-\(k\) trace polynomial.

\begin{remark}
For fixed \(\varepsilon>0\), the second term is much smaller than
\(\exp(-c d^2)\).  We keep the two terms separated because the proof naturally
produces them from different mechanisms: probabilistic construction of the
good set and spherical concentration on the good set.
\end{remark}

\section{Proof architecture}

The proof of Theorem~\ref{thm:main-concentration} has four conceptual steps.

\begin{enumerate}[leftmargin=*]
\item \emph{Localize the concentration inequality.}
  Levy concentration is applied not directly to \(f_k\), but to a Lipschitz
  extension of \(f_k\) from a high-probability good set.  This produces a tail
  estimate with two errors: the spherical Levy tail and the probability of the
  bad set.

\item \emph{Build the good set.}
  On the good set, the operator norm of \((XX^*)^\Gamma\) and the relevant
  diagonal gamma quantities are bounded at the correct \(d^{-2}\) scale.  The
  complement of this good set is exponentially small.

\item \emph{Prove the deterministic Lipschitz bound on the good set.}
  The trace-power telescoping identity reduces the variation of \(f_k\) to a
  trace-norm perturbation bound.  Hilbert--Schmidt ideal inequalities and the
  operator-norm control on the good set yield the desired local Lipschitz
  constant.

\item \emph{Replace the median by the expectation.}
  The localized concentration estimate first gives concentration around a
  median.  Since \(f_k\) is uniformly bounded, the tail estimate can be
  integrated to show that the median and the mean differ by a smaller error.
  This gives the final concentration inequality around \(\E f_k\).
\end{enumerate}

\subsection{Why the local Lipschitz constant has the right scale}

The deterministic perturbation bound starts from
\[
  \abs{\Tr(A^k)-\Tr(B^k)}
  \le
  k\max(\norm A_\infty,\norm B_\infty)^{k-1}\norm{A-B}_1.
\]
Applied to
\[
  A=(XX^*)^\Gamma,\qquad B=(YY^*)^\Gamma,
\]
the trace-norm difference is controlled by
\[
  \norm{A-B}_1
  \lesssim
  d\,\max(\norm X_\infty,\norm Y_\infty)\norm{X-Y}_2.
\]
On the good set one has
\[
  \norm{(XX^*)^\Gamma}_\infty \lesssim d^{-2},
  \qquad
  \norm X_\infty \lesssim d^{-1}.
\]
Consequently
\[
  \abs{f_k(X)-f_k(Y)}
  \lesssim
  k d^{-2(k-1)}\norm{X-Y}_2
  \qquad (X,Y\in\Omega).
\]
Thus the relevant Lipschitz constant is
\[
  L_{k,d}\asymp k d^{-2k+2}.
\]
Since the sphere has real dimension \(n\asymp d^4\), Levy's lemma gives
\[
  \exp\left(-c\,\frac{n t^2}{L_{k,d}^2}\right)
  =
  \exp\left(-c\,\frac{d^4\varepsilon^2}{k^2}\right)
\]
when \(t=\varepsilon d^{-(2k-2)}\).  This is the scale appearing in
Theorem~\ref{thm:main-concentration}.

\subsection{The good set}

The good set \(\Omega\) is the intersection of two events.  The first controls
the operator norm of \(X\), and follows from standard Gaussian operator-norm
estimates after writing \(X=G/\norm G_2\).  The second controls
\((XX^*)^\Gamma\).  For this, one decomposes the partially transposed Wishart
matrix into its diagonal and off-diagonal parts.  The diagonal part is
controlled by gamma concentration, while the off-diagonal part is controlled
by Aubrun's moment estimate and Markov's inequality applied to a high even
moment.

The output is the estimate
\[
  \Prob(\Omega^c)\le C e^{-c d^2}.
\]
This probability is not negligible compared with the final bound; it is
exactly the first term in Theorem~\ref{thm:main-concentration}.

\subsection{Median and expectation}

The localized Levy argument first gives concentration around a median.  To
pass to the expectation, we use the uniform bound \(\abs{f_k(X)}\le1\) for
\(k\ge2\), while \(f_1(X)=1\).  Integrating the median tail only up to the
bounded range gives
\[
  \abs{\E f_k-M_k}
  \le
  C e^{-c d^2}+C_{k,\lambda}k d^{-2k}.
\]
This error is smaller than the target scale
\(\varepsilon d^{-(2k-2)}\) except in the regime where the stated probability
bound is already trivial after increasing the constant \(C\).  This completes
the reduction from median concentration to mean concentration.

\subsection{Detailed proof}

The rest of the article gives the detailed estimates in the order just
described.
\end{revderived}

\appendix
\section{Concentration via local Lipschitz estimates}
\label{app:second_approach}

\makeatletter
\@ifundefined{revderived}{\newenvironment{revderived}{}{}}{}
\@ifundefined{revlean}{\newenvironment{revlean}{}{}}{}
\@ifundefined{revnewmath}{\newenvironment{revnewmath}{}{}}{}
\@ifundefined{revlatest}{\newenvironment{revlatest}{}{}}{}
\@ifundefined{revrecent}{\newenvironment{revrecent}{}{}}{}
\@ifundefined{revearlier}{\newenvironment{revearlier}{}{}}{}
\makeatother

We prove concentration for the normalized moments
\[
  f_k(X)=\Tr\left(\left((XX^*)^\Gamma\right)^k\right),
  \qquad
  X\in S_{HS}(d^2,s):=\{X\in M_{d^2,s}(\C):\norm X_2=1\}.
\]
The sphere \(S_{HS}(d^2,s)\) is equipped with its uniform probability
measure, i.e. Euclidean surface measure divided by its surface area.

We use the following norm conventions.  For \(X\in M_{m,n}(\C)\),
\[
  \norm X_2=\bigl(\Tr(X^*X)\bigr)^{1/2}
  =
  \left(\sum_{i=1}^m\sum_{\alpha=1}^n |X_{i\alpha}|^2\right)^{1/2}
\]
is the Hilbert--Schmidt norm, and
\[
  \norm X_\infty
  =
  \sup_{\norm v_{\ell_2^n}=1}\norm{Xv}_{\ell_2^m}
\]
is the operator norm.  For \(A\in M_N(\C)\),
\[
  \norm A_1=\Tr|A|,\qquad |A|=(A^*A)^{1/2},
\]
is the trace norm.  We regard \(M_{d^2}(\C)\) as
\(M_d(\C)\otimes M_d(\C)\), and
\(\Gamma=\operatorname{Id}\otimes T\) denotes partial
transpose on the second tensor factor.

The only external estimates used below are Levy's spherical concentration
inequality \cite[Corollary~5.17]{AubrunSzarek2017} and Aubrun's
off-diagonal partially-transposed Wishart moment estimate
\cite[Lemma~3.4]{Aubrun12}; all other reductions are written out.

\begin{theorem}[Localized concentration of normalized moments]
\label{thm:appendixB-local-concentration}
For every fixed \(k\ge1\) and every \(\lambda>0\), there are constants
\(c,C>0\), depending only on \(k\) and \(\lambda\), such that for all
sufficiently large \(d\), all integers \(s\) satisfying
\[
  \frac{\lambda}{2}d^2\le s\le 2\lambda d^2,
\]
and all \(\varepsilon>0\),
\[
  \Prob\left\{
    \left|
      \Tr\left(\left((XX^*)^\Gamma\right)^k\right)
      -
      \E\Tr\left(\left((XX^*)^\Gamma\right)^k\right)
    \right|
    \ge
    \varepsilon d^{-(2k-2)}
  \right\}
  \le
  C e^{-c d^2}
  +
  C\exp\left(-c\frac{d^4\varepsilon^2}{k^2}\right),
\]
where \(X\) is uniform on \(S_{HS}(d^2,s)\).
\end{theorem}

\subsection{A local form of Levy's lemma}

Let \(\sigma_{n-1}\) denote the uniform probability measure on the Euclidean
unit sphere \(S^{n-1}\subset\R^n\), namely Euclidean surface measure divided
by the surface area of \(S^{n-1}\).  Distances on \(S^{n-1}\) are measured
with the ambient Euclidean norm unless explicitly stated otherwise.

\begin{lemma}[Spherical Levy inequality]
\label{lem:appendixB-spherical-levy}
Assume \(n\ge3\).  Let \(F:S^{n-1}\to\R\) be \(L\)-Lipschitz for the
ambient Euclidean distance, with \(L>0\), and let \(M_F\) be a median of
\(F\) with respect to \(\sigma_{n-1}\).  Then for every \(t>0\),
\[
  \sigma_{n-1}\{x\in S^{n-1}: |F(x)-M_F|> t\}
  \le \exp\left(-\frac{n t^2}{2L^2}\right)
\].
\end{lemma}

\begin{proof}
Aubrun--Szarek's Levy lemma, \cite[Corollary~5.17]{AubrunSzarek2017},
states the following stronger median estimate for the geodesic distance
\(g\) on \(S^{n-1}\): if \(n>2\), if \(F\) is \(L\)-Lipschitz for \(g\),
and if \(M_F\) is a median of \(F\), then, for every \(t>0\),
\[
  \sigma_{n-1}\{x: |F(x)-M_F|>t\}
  \le
  \exp\left(-\frac{n t^2}{2L^2}\right).
\]
The ambient Euclidean distance is dominated by the geodesic distance:
\[
  \norm{x-y}_2\le g(x,y)\qquad(x,y\in S^{n-1}).
\]
Thus an \(L\)-Lipschitz function for the Euclidean distance is also
\(L\)-Lipschitz for \(g\), so the displayed estimate is exactly the stated
estimate.
\end{proof}

The following elementary localization is the form needed below; all
probabilities in the statement are with respect to \(\sigma_{n-1}\).

\begin{lemma}[Localized Levy lemma]
\label{lem:appendixB-local-levy}
Assume \(n\ge3\).  Let \(\Omega\subset S^{n-1}\) be a nonempty measurable
set, and let
\(f:S^{n-1}\to\R\) be measurable.  Assume that \(f\) is
\(L\)-Lipschitz on \(\Omega\), with \(L>0\).  If \(M_f\) is a median of
\(f\), then for every \(t>0\),
\[
  \Prob\{|f-M_f|>t\}
  \le
  2\Prob(\Omega^c)
  +
  2\exp\left(-\frac{n t^2}{8L^2}\right).
\]
\end{lemma}

\begin{proof}
Define the McShane extension of \(f|_\Omega\) by
\[
  F(x)=\inf_{y\in\Omega}\bigl(f(y)+L\norm{x-y}_2\bigr),
  \qquad x\in S^{n-1}.
\]
Then \(F=f\) on \(\Omega\) and \(F\) is \(L\)-Lipschitz on the whole sphere.
Indeed, for any \(x,x'\) and any \(y\in\Omega\),
\[
  f(y)+L\norm{x-y}_2
  \le f(y)+L\norm{x'-y}_2+L\norm{x-x'}_2,
\]
and taking the infimum over \(y\) gives
\(F(x)\le F(x')+L\norm{x-x'}_2\); reversing \(x,x'\) gives the Lipschitz
bound.  If \(x\in\Omega\), the choice \(y=x\) gives \(F(x)\le f(x)\), while
the Lipschitz property of \(f|_\Omega\) gives
\(f(x)\le f(y)+L\norm{x-y}_2\) for every \(y\in\Omega\), hence
\(f(x)\le F(x)\).

Let \(M_F\) be a median of \(F\).  Since \(F=f\) on \(\Omega\),
\[
  \Prob\{|f-M_F|>u\}
  \le
  \Prob(\Omega^c)+
  \exp\left(-\frac{n u^2}{2L^2}\right).
\]
Here the exponential term is exactly Lemma~\ref{lem:appendixB-spherical-levy}
applied to the global \(L\)-Lipschitz function \(F\); the extra
\(\Prob(\Omega^c)\) accounts for the points where \(F\) and \(f\) need not
agree.
The median replacement estimate
\[
  \Prob\{|f-M_f|>t\}
  \le
  2\Prob\{|f-M_F|>t/2\}
\]
is elementary: if \(|M_f-M_F|\le t/2\), the event on the left is contained
in \(\{|f-M_F|>t/2\}\); if \(|M_f-M_F|>t/2\), then
\(\Prob\{|f-M_F|>t/2\}\ge1/2\) by the definition of a median.  Combining the
two estimates gives the claim.
\end{proof}

\subsection{Deterministic Lipschitz estimate}

\begin{lemma}[Trace-power perturbation]
\label{lem:appendixB-trace-power-perturbation}
Let \(N,k\ge1\).  For every \(A,B\in M_N(\C)\),
\[
  |\Tr(A^k)-\Tr(B^k)|
  \le
  k\max(\norm A_\infty,\norm B_\infty)^{k-1}\norm{A-B}_1 .
\]
\end{lemma}

\begin{proof}
The telescoping identity gives
\[
  A^k-B^k
  =
  \sum_{j=0}^{k-1}A^j(A-B)B^{k-1-j}.
\]
For every matrix \(M\), \(|\Tr M|\le\norm M_1\), and the trace norm satisfies
\[
  \norm{RST}_1\le\norm R_\infty\norm S_1\norm T_\infty .
\]
Therefore
\[
\begin{aligned}
|\Tr(A^k)-\Tr(B^k)|
&\le
\sum_{j=0}^{k-1}
\norm{A^j(A-B)B^{k-1-j}}_1       \\
&\le
\sum_{j=0}^{k-1}
\norm A_\infty^j\norm{A-B}_1\norm B_\infty^{k-1-j}                 \\
&\le
k\max(\norm A_\infty,\norm B_\infty)^{k-1}\norm{A-B}_1 .
\end{aligned}
\]
\end{proof}

\begin{lemma}[Partial-transpose moment map Lipschitz bound]
\label{lem:appendixB-moment-lipschitz}
Let \(d,s,k\ge1\).  For \(X,Y\in M_{d^2,s}(\C)\), put
\[
  A_X=(XX^*)^\Gamma,\qquad A_Y=(YY^*)^\Gamma.
\]
Then
\[
\begin{aligned}
&\left|
  \Tr(A_X^k)-\Tr(A_Y^k)
\right|                                                     \\
&\qquad\le
2kd\,
\max(\norm{A_X}_\infty,\norm{A_Y}_\infty)^{k-1}
\max(\norm X_\infty,\norm Y_\infty)\,
\norm{X-Y}_2 .
\end{aligned}
\]
\end{lemma}

\begin{proof}
By Lemma~\ref{lem:appendixB-trace-power-perturbation},
\[
\left|\Tr(A_X^k)-\Tr(A_Y^k)\right|
\le
k\max(\norm{A_X}_\infty,\norm{A_Y}_\infty)^{k-1}
\norm{A_X-A_Y}_1 .
\]
Since \(A_X-A_Y=(XX^*-YY^*)^\Gamma\), partial transpose is an isometry for
the Hilbert--Schmidt norm, and
\(\norm M_1\le d\norm M_2\) for \(M\in M_{d^2}(\C)\), we get
\[
  \norm{A_X-A_Y}_1
  \le d\norm{XX^*-YY^*}_2.
\]
Finally,
\[
  XX^*-YY^*=(X-Y)X^*+Y(X-Y)^*,
\]
and the Hilbert--Schmidt ideal inequality gives
\[
  \norm{XX^*-YY^*}_2
  \le
  \norm{X-Y}_2\norm X_\infty+\norm Y_\infty\norm{X-Y}_2
  \le
  2\max(\norm X_\infty,\norm Y_\infty)\norm{X-Y}_2.
\]
Combining the estimates proves the lemma.
\end{proof}

\subsection{Operator-norm good sets}

\begin{lemma}[Gaussian polar decomposition]
\label{lem:appendixB-gaussian-polar}
Let \(m,n\ge1\), and let \(G\in M_{m,n}(\C)\) have independent standard
complex Gaussian entries.  Put \(R=\norm G_2\) and \(X=G/R\), on the
almost-sure event \(\{R>0\}\).  Then \(X\) is distributed according to the
uniform probability measure on
\[
  S_{HS}(m,n):=\{Y\in M_{m,n}(\C):\norm Y_2=1\},
\]
and \(X\) is independent of \(R\).  Moreover
\[
  R^2=\sum_{i,\alpha}|G_{i\alpha}|^2
\]
has Gamma distribution with shape \(mn\) and scale \(1\).
\end{lemma}

\begin{proof}
Identify \(M_{m,n}(\C)\) with \(\C^{mn}\), equipped with the Euclidean norm
\(\norm{\cdot}_2\).  The density of \(G\) with respect to Lebesgue measure is
\[
  z\longmapsto \pi^{-mn}\exp(-\norm z_2^2).
\]
The origin has Gaussian measure zero, so \(R>0\) almost surely.  In polar
coordinates \(z=ru\), with \(r>0\) and \(u\in S_{HS}(m,n)\), Lebesgue measure
has the form
\[
  dz=c_{m,n}\, r^{2mn-1}\,dr\,d\sigma(u),
\]
where \(\sigma\) is the uniform probability measure on \(S_{HS}(m,n)\) and
\(c_{m,n}\) is the surface area constant.  Hence the joint density of
\((R,X)\) is proportional to
\[
  e^{-r^2}r^{2mn-1}\,dr\,d\sigma(u).
\]
This factorizes into a function of \(r\) times \(\sigma(du)\), proving both
the uniformity of \(X\) and the independence of \(X\) and \(R\).  Finally,
the change of variable \(y=r^2\) gives density proportional to
\[
  y^{mn-1}e^{-y}\,dy
\]
for \(R^2\), i.e. the Gamma distribution with shape \(mn\) and scale \(1\).
\end{proof}

\begin{lemma}[Gaussian operator norm]
\label{lem:appendixB-gaussian-op}
Let \(m,n\ge1\), and let \(H\in M_{m,n}(\C)\) have independent standard
complex Gaussian entries.  There is a universal constant \(C\) such that
\[
  \E\norm H_\infty\le C(\sqrt m+\sqrt n).
\]
\end{lemma}

\begin{proof}
For unit vectors \(u\in\C^m\) and \(v\in\C^n\),
\[
  \langle u,Hv\rangle
  =
  \sum_{i,\alpha}\overline{u_i}H_{i\alpha}v_\alpha
\]
is a standard complex Gaussian, so
\(\Prob\{|\langle u,Hv\rangle|\ge r\}=e^{-r^2}\).
Choose \(1/4\)-nets \(N_1,N_2\) in the unit spheres of \(\C^m\) and
\(\C^n\).  The elementary volume argument gives
\[
  |N_1|\le9^{2m},\qquad |N_2|\le9^{2n}.
\]
If
\[
  M=\max_{u\in N_1,\ v\in N_2}|\langle u,Hv\rangle|,
\]
then \(\norm H_\infty\le2M\).  To see this, approximate arbitrary unit
\(u,v\) by \(u_0,v_0\) in the nets and use
\[
  |\langle u,Hv\rangle|
  \le
  |\langle u_0,Hv_0\rangle|
  +\frac14\norm H_\infty+\frac14\norm H_\infty .
\]
By the union bound,
\[
  \Prob\{M\ge r\}\le
  \exp\{2(m+n)\log9-r^2\}.
\]
Integrating this tail gives
\(\E M\le C'(\sqrt m+\sqrt n)\), and hence
\(\E\norm H_\infty\le C(\sqrt m+\sqrt n)\).
\end{proof}

\begin{revlatest}
\begin{lemma}[Median bound for nonnegative random variables]
\label{lem:appendixB-nonnegative-median-bound}
Let \(Z\ge0\) be an integrable random variable and let \(M\in\R\) satisfy
\[
  \Prob\{Z\ge M\}\ge \frac12 .
\]
Then
\[
  M\le 2\E Z .
\]
In particular, every median of a nonnegative random variable is bounded above
by twice its expectation.
\end{lemma}

\begin{proof}
If \(M\le0\), the claim is immediate.  Otherwise \(M>0\), and
\[
  \E Z
  \ge
  \E\bigl(Z\mathbf 1_{\{Z\ge M\}}\bigr)
  \ge
  M\Prob\{Z\ge M\}
  \ge
  \frac M2 .
\]
\end{proof}
\end{revlatest}

\begin{proposition}[Partial-transpose Wishart expectation in the balanced window]
\label{prop:appendixB-wishart-expectation}
For every \(\lambda>0\), there are constants \(C_\lambda>0\) and
\(d_0=d_0(\lambda)\) such that the following holds.  Let \(d\ge d_0\), let
\(s\) be an integer satisfying
\[
  \frac{\lambda}{2}d^2\le s\le 2\lambda d^2 .
\]
Let \(G\in M_{d^2,s}(\C)\) have independent standard complex Gaussian
entries and put \(W=s^{-1}GG^*\).  Then
\[
  \E\norm{W^\Gamma}_\infty\le C_\lambda .
\]
\end{proposition}

\begin{proof}
Write
\[
  W^\Gamma=\diag(W^\Gamma)+Z,\qquad Z=W^\Gamma-\diag(W^\Gamma).
\]
The diagonal is unchanged by partial transpose:
\[
  (W^\Gamma)_{ii}=W_{ii}
  =
  \frac1s\sum_{\alpha=1}^s |G_{i\alpha}|^2.
\]
Thus \(\norm{\diag(W^\Gamma)}_\infty\) is the maximum of \(d^2\) averages of
\(s\) independent mean-one exponentials.  We record the Chernoff estimate
needed here.  If \(V_i=s^{-1}\sum_{\alpha=1}^s E_{i\alpha}\), with the
\(E_{i\alpha}\)'s independent exponentials of mean one, then, for every
\(u>0\),
\[
  \Prob\left\{
    V_i\ge 1+2\sqrt{u/s}+2u/s
  \right\}\le e^{-u}.
\]
Indeed, for \(\theta\in(0,1)\),
\[
  \Prob\{V_i\ge 1+\delta\}
  \le
  e^{-\theta s(1+\delta)}(1-\theta)^{-s}.
\]
Taking \(\theta=\delta/(1+\delta)\) gives
\[
  \Prob\{V_i\ge 1+\delta\}
  \le
  \exp\{-s(\delta-\log(1+\delta))\}.
\]
With \(z=\sqrt{u/s}\) and \(\delta=2z+2z^2\), the elementary inequality
\[
  \log(1+2z+2z^2)\le 2z+z^2
\]
implies \(\delta-\log(1+\delta)\ge z^2=u/s\), which proves the tail bound.

Apply this with \(u=v+2\log d\) and take a union bound over the \(d^2\)
diagonal entries.  For every \(v\ge0\),
\[
  \Prob\left\{
    \max_i V_i
    \ge
    1+2\sqrt{\frac{v+2\log d}{s}}
      +2\frac{v+2\log d}{s}
  \right\}
  \le e^{-v}.
\]
Let
\[
  a(v)=
  1+2\sqrt{\frac{v+2\log d}{s}}
    +2\frac{v+2\log d}{s}.
\]
The function \(a\) is increasing, and
\[
  a'(v)=
  \frac1{\sqrt{s}\sqrt{v+2\log d}}+\frac2s.
\]
By the layer-cake formula and the preceding tail bound,
\[
\begin{aligned}
  \E\max_{1\le i\le d^2}V_i
  &\le
  a(0)+\int_0^\infty e^{-v}a'(v)\,dv                         \\
  &\le
  1+2\sqrt{\frac{2\log d}{s}}
    +4\frac{\log d}{s}
    +\frac1{\sqrt{s}}\int_0^\infty \frac{e^{-v}}{\sqrt v}\,dv
    +\frac2s                                                    \\
  &\le
  1+C\sqrt{\frac{\log d}{s}}
    +C\frac{\log d}{s}
\end{aligned}
\]
for \(d\ge2\).
Therefore
\[
  \E\norm{\diag(W^\Gamma)}_\infty
  \le
  1+C_\lambda\frac{\sqrt{\log d}}{d}.
\]

For the off-diagonal part we use Aubrun's moment estimate,
\cite[Lemma~3.4]{Aubrun12}: for every even \(m\ge2\) there is a polynomial
\(Q\), independent of \(d\) and \(s\), such that
\[
  \E\Tr(Z^m)
  \le
  \left(\frac2s\right)^m
  (d+Q(m))^{m+2}(\sqrt s+Q(m))^m.
\]
Since \(Z\) is Hermitian and \(m\) is even,
\[
  \E\norm Z_\infty
  \le
  \bigl(\E\Tr(Z^m)\bigr)^{1/m}.
\]
Choose even \(m=m(d)\to\infty\) slowly enough that
\[
  Q(m)=o(d),\qquad \frac{\log d}{m}\to0.
\]
Taking the \(m\)-th root of Aubrun's estimate gives
\[
  \E\norm Z_\infty
  \le
  \frac2s(d+Q(m))^{1+2/m}(\sqrt s+Q(m)),
\]
which is bounded uniformly in \(d\) because
\(\lambda d^2/2\le s\le2\lambda d^2\).  The triangle inequality gives the
result.
\end{proof}

\begin{revlatest}
\begin{lemma}[Radial transfer from Wishart to the sphere]
\label{lem:appendixB-radial-transfer-pt-operator}
Let \(G\in M_{d^2,s}(\C)\) have independent standard complex Gaussian
entries, and use its polar decomposition in the real Hilbert space
\(M_{d^2,s}(\C)\) with the Hilbert--Schmidt norm.  Thus, outside the null
event \(\{G=0\}\), set
\[
  R=\norm G_2,\qquad X=\frac GR,\qquad W=\frac1s GG^*
\]
so that \(G=RX\), with \(X\in S_{HS}(d^2,s)\).  By
Lemma~\ref{lem:appendixB-gaussian-polar}, \(X\) is uniform on the
Hilbert--Schmidt sphere and independent of \(R\).  Then
\[
  \E\norm{(XX^*)^\Gamma}_\infty
  =
  \frac1{d^2}\E\norm{W^\Gamma}_\infty .
\]
\end{lemma}

\begin{proof}
By Lemma~\ref{lem:appendixB-gaussian-polar}, \(\E R^2=d^2s\).  Since
\(G=RX\), we have
\[
  W^\Gamma
  =
  \frac1s(GG^*)^\Gamma
  =
  \frac{R^2}{s}(XX^*)^\Gamma .
\]
Taking operator norms and expectations gives
\[
  \E\norm{W^\Gamma}_\infty
  =
  \frac1s\E\left(R^2\norm{(XX^*)^\Gamma}_\infty\right)
  =
  \frac{\E R^2}{s}\E\norm{(XX^*)^\Gamma}_\infty
  =
  d^2\,\E\norm{(XX^*)^\Gamma}_\infty,
\]
where the middle equality uses independence of \(R\) and \(X\).  Rearranging
proves the lemma.
\end{proof}
\end{revlatest}

\begin{proposition}[Good set in the balanced window]
\label{prop:appendixB-good-set}
For every \(\lambda>0\), there are constants \(a,b,c,C>0\) and
\(d_0=d_0(\lambda)\) such that the following holds.  Let \(d\ge d_0\), let
\(s\) be an integer satisfying \(\lambda d^2/2\le s\le2\lambda d^2\), and
let \(X\) be uniform on \(S_{HS}(d^2,s)\).  Then the set
\[
  \Omega
  =
  \left\{
    X:\norm X_\infty\le\frac a d,\quad
    \norm{(XX^*)^\Gamma}_\infty\le\frac b{d^2}
  \right\}
\]
satisfies
\[
  \Prob(\Omega^c)\le C e^{-c d^2}.
\]
\end{proposition}

\begin{proof}
Let \(G\in M_{d^2,s}(\C)\) have independent standard complex Gaussian
entries.  On the almost-sure event \(G\ne0\), put
\[
  R=\norm G_2,\qquad X=\frac GR .
\]
The polar law in Lemma~\ref{lem:appendixB-gaussian-polar} says precisely that
this angular variable \(X\) has the uniform surface law on \(S_{HS}(d^2,s)\)
and is independent of the radial variable \(R\).
First consider \(h(X)=\norm X_\infty\).  It is \(1\)-Lipschitz on the
Hilbert--Schmidt sphere.  By independence of \(R\) and \(X\),
\[
  \E\norm G_\infty=\E R\,\E\norm X_\infty.
\]
Since \(R^2\) has Gamma distribution with shape \(N=d^2s\) and scale \(1\),
\[
  \E R^2=N,\qquad \E R^4=N(N+1).
\]
The Paley--Zygmund inequality applied to \(R^2\) gives
\[
  \Prob\{R^2\ge N/2\}
  \ge
  \frac{(N/2)^2}{N(N+1)}
  \ge \frac18 .
\]
Consequently
\[
  \E R\ge c\sqrt{d^2s}.
\]
Together with Lemma~\ref{lem:appendixB-gaussian-op}, applied with
\(m=d^2\) and \(n=s\), this gives
\[
  \E\norm X_\infty\le \frac{C_\lambda}{d}.
\]
\begin{revlatest}
By Lemma~\ref{lem:appendixB-nonnegative-median-bound}, any median \(M_h\)
satisfies \(M_h\le2\E h\).  Levy's lemma in the form of
Lemma~\ref{lem:appendixB-spherical-levy}, applied on the real sphere of
dimension \(2d^2s-1\), gives, with deviation \(C_\lambda/d\),
\end{revlatest}
\[
  \Prob\left\{\norm X_\infty>\frac a d\right\}\le e^{-c d^2}.
\]

Now set \(g(X)=\norm{(XX^*)^\Gamma}_\infty\).  On the event
\(\Omega_1=\{\norm X_\infty\le a/d\}\), the map \(g\) is
\(C_\lambda/d\)-Lipschitz, because
\[
\begin{aligned}
 |g(X)-g(Y)|
 &\le
 \norm{(XX^*)^\Gamma-(YY^*)^\Gamma}_\infty        \\
 &\le
 \norm{XX^*-YY^*}_2
 \le
 (\norm X_\infty+\norm Y_\infty)\norm{X-Y}_2 .
\end{aligned}
\]
\begin{revlatest}
Furthermore, Lemma~\ref{lem:appendixB-radial-transfer-pt-operator} and
Proposition~\ref{prop:appendixB-wishart-expectation} give
\end{revlatest}
\[
  \E g(X)
  \le
  \frac{C_\lambda}{d^2}.
\]
\begin{revlatest}
Again Lemma~\ref{lem:appendixB-nonnegative-median-bound} bounds any median
\(M_g\) by \(2\E g\).  Applying the localized Levy lemma to \(g\) on
\(\Omega_1\), with deviation \(C_\lambda/d^2\), gives
\end{revlatest}
\[
  \Prob\left\{
    \norm{(XX^*)^\Gamma}_\infty>\frac b{d^2}
  \right\}
  \le
  C e^{-c d^2}.
\]
The union bound proves the proposition.
\end{proof}

\subsection{Proof of Theorem~\ref{thm:appendixB-local-concentration}}

On the good set \(\Omega\) from Proposition~\ref{prop:appendixB-good-set},
Lemma~\ref{lem:appendixB-moment-lipschitz} gives
\[
  |f_k(X)-f_k(Y)|
  \le
  \frac{C_{k,\lambda}k}{d^{2k-2}}\norm{X-Y}_2
  \qquad (X,Y\in\Omega).
\]
The real dimension of the Hilbert--Schmidt sphere is \(2d^2s-1\), which is
of order \(d^4\) and is at least \(2\) for all sufficiently large \(d\).
Applying Lemma~\ref{lem:appendixB-local-levy} with
\[
  L=\frac{C_{k,\lambda}k}{d^{2k-2}},
  \qquad
  t=\frac12\varepsilon d^{-(2k-2)},
\]
and using
\(\{|f_k-M_k|\ge \varepsilon d^{-(2k-2)}\}
  \subset
  \{|f_k-M_k|>\frac12\varepsilon d^{-(2k-2)}\}\),
we obtain, for every median \(M_k\) of \(f_k\),
\[
  \Prob\left\{
    |f_k(X)-M_k|\ge \varepsilon d^{-(2k-2)}
  \right\}
  \le
  C e^{-c d^2}
  +
  4\exp\left(-c\frac{d^4\varepsilon^2}{k^2}\right).
\]

It remains to replace the median by the expectation.  Since partial
transpose is an isometry for the Hilbert--Schmidt norm,
\[
  \norm{(XX^*)^\Gamma}_2=\norm{XX^*}_2.
\]
Moreover
\[
  \norm{XX^*}_2^2=\Tr((X^*X)^2)
  \le \Tr(X^*X)^2=\norm X_2^4=1.
\]
Thus, if \((\alpha_j)\) are the eigenvalues of \((XX^*)^\Gamma\), then
\(\sum_j|\alpha_j|^2\le1\).  Therefore \(|f_k(X)|\le1\) for \(k\ge2\),
while \(f_1(X)=1\).  From now on choose a median \(M_k\in[-1,1]\); this is
possible because \(f_k\) takes values in \([-1,1]\).  Hence
\[
  |f_k(X)-M_k|\le2
\]
for every \(X\).  This boundedness is not used to integrate a constant tail
on \([0,\infty)\).  It is used to truncate the layer-cake identity:
\[
\begin{aligned}
  |\E f_k-M_k|
  &\le \E |f_k-M_k|                                      \\
  &=\int_0^2
      \Prob\{|f_k(X)-M_k|\ge u\}\,du .
\end{aligned}
\]
For \(0<u\le2\), the localized Levy estimate applied with this value of
\(u\), together with \(n=2d^2s\asymp_\lambda d^4\) and
\(L=C_{k,\lambda}k d^{-(2k-2)}\), gives
\[
  \Prob\{|f_k(X)-M_k|\ge u\}
  \le
  C e^{-c d^2}
  +
  4\exp\left(-c\,\frac{d^{4k}u^2}{k^2}\right).
\]
Therefore
\[
\begin{aligned}
  |\E f_k-M_k|
  &\le
  2C e^{-c d^2}
  +
  4\int_0^\infty
    \exp\left(-c\,\frac{d^{4k}u^2}{k^2}\right)\,du        \\
  &\le
  C e^{-c d^2}+\frac{C_{k,\lambda}k}{d^{2k}} .
\end{aligned}
\]
If the right hand side in the last display is at most
\(\varepsilon d^{-(2k-2)}/2\), then
\[
  \{|f_k-\E f_k|\ge\varepsilon d^{-(2k-2)}\}
  \subset
  \{|f_k-M_k|\ge\tfrac12\varepsilon d^{-(2k-2)}\},
\]
and the preceding median estimate gives the result after changing constants.
If instead
\[
  |\E f_k-M_k|>\frac12\varepsilon d^{-(2k-2)},
\]
then the displayed bound on \(|\E f_k-M_k|\) forces
\(\varepsilon\le C_{k,\lambda}d^{-2}\), up to changing constants and ignoring
the exponentially small term.  In that regime
\[
  \exp\left(-c\frac{d^4\varepsilon^2}{k^2}\right)
\]
is bounded below by a positive constant depending only on \(k,\lambda\), so
the right hand side of the theorem can be made at least \(1\) by increasing
\(C\).  The probability bound is then trivial.  This proves the theorem.


% Standalone-only bibliography for compiling this wrapper outside the article.
% Do not paste this block into a host article that already uses a bibliography.
\begin{thebibliography}{2}
\bibitem{Aubrun12}
G.~Aubrun,
\emph{Partial transposition of random states and non-centered semicircular
distributions},
Random Matrices: Theory Appl. \textbf{1} (2012), 1250001.

\bibitem{AubrunSzarek2017}
G.~Aubrun and S.~J. Szarek,
\emph{Alice and Bob Meet Banach: The Interface of Asymptotic Geometric
Analysis and Quantum Information Theory},
Mathematical Surveys and Monographs, vol.~223, American Mathematical
Society, 2017.
\end{thebibliography}

\end{document}
